\documentclass[10pt,a4paper]{amsart}
\usepackage[margin=19mm]{geometry}
\usepackage{amsmath,amssymb,mathtools}
\usepackage[scr=rsfs,cal=euler]{mathalfa}
\usepackage{xfrac}
\usepackage[unicode,hidelinks,bookmarks=false]{hyperref}
\newcommand{\ie}{i.e.\ }
\newcommand{\cf}{cf.\ }
\newcommand{\resp}{respectively\ }
\newcommand{\Subsection}{\subsection}
\providecommand{\nobreakdash}{\nobreak\hbox{-}}
\setlength{\emergencystretch}{2em}
\multlinegap0pt
\usepackage{bm}
\usepackage{bbm}
\usepackage{dsfont}
\usepackage{xspace}
\usepackage{cancel}
\usepackage{mathtools}
\usepackage{amsthm}
\makeatletter
\@ifundefined{theorem}{\newtheorem{theorem}{Theorem}}{}
\@ifundefined{lemma}{\newtheorem{lemma}[theorem]{Lemma}}{}
\makeatother
\title[Equivariant cohomology of even-dimensional complex quadrics ]{Equivariant cohomology of even-dimensional complex quadrics from a combinatorial point of view}
\author{Shintaro Kuroki}
\date{Source: arXiv:2305.11332v3 (CC BY 4.0)}
\begin{document}
\maketitle

\noindent\emph{Source and licence.} Equivariant cohomology of even-dimensional complex quadrics from a combinatorial point of view. Version arXiv:2305.11332v3 (CC BY 4.0), distributed under the Creative Commons Attribution 4.0 International licence, as recorded in the arXiv licence metadata for this exact version and shown on the abstract page https://arxiv.org/abs/2305.11332v3. Full rights notice: licenses/arxiv-2305.11332-CC-BY-4.0.txt.\par\smallskip

\noindent\textit{Editorial scope.} Counted unit: the Lemma whose source label is \texttt{injectivity} in the article (source0.tex lines 2334--2339, source label \texttt{injectivity}), delivered together with the article's own complete original proof of it (source0.tex lines 2463--2595, one \texttt{proof} environment of 133 lines). No mathematics is written, repaired or re-derived: statement and proof are the article's own text line for line. The proof is detached from the statement in the source: the article states the result at lines 2334--2339 and prints its proof at lines 2463--2595, and the proof environment's own header names the statement, so the association is the article's own. Required context retained uncounted: algebra-structure (lines 700--703); surjectivity (lines 2191--2196); linear-relation (lines 2311--2316). Every definition, display and label of those lines is retained unchanged. Omitted from this package and not redistributed: 1--699, 704--2190, 2197--2310, 2317--2333, 2340--2462, 2596--3459. Nothing inside a retained excerpt is omitted or altered: every retained body is delivered exactly as the article's lines are written, so no omission receipt of any kind accompanies this package. Citations: the retained excerpts carry no citation command at all, so no bibliography entry of the article is transcribed and no reference is redistributed. Reference adaptation: none was needed and none was made; every reference inside the delivered unit resolves inside it, so no unresolved reference or citation is produced. Printed slips kept literal: no printed word, symbol, hypothesis or number of the counted statement or of its proof was altered. Layout and class: the article's own class is \texttt{amsart} with packages including a4wide, color, amssymb, amsmath, mathrsfs, float, graphicx, tikz, tikz-cd, epsfig, xy, epstopdf, framed; the wrapper is the lane's pinned \texttt{amsart} editorial wrapper at 10pt on A4, which restates the theorem-like environments the retained text needs and adds the article's own macro definitions that the retained text needs (none), with extra wrapper packages (bm, bbm, dsfont, xspace, cancel, mathtools). The delivered numbering is the unit's own. Assets: no retained span contains a figure environment, a graphics-inclusion command, a PDF-page inclusion, a table or a TikZ picture; no image file is copied into this package, the package declares no asset dependency and no image file is redistributed with it. Licence: version arXiv:2305.11332v3 of this article is distributed under the Creative Commons Attribution 4.0 International licence, as recorded in the arXiv licence metadata for this exact version and shown on the abstract page https://arxiv.org/abs/2305.11332v3. A full rights notice is delivered at licenses/arxiv-2305.11332-CC-BY-4.0.txt. Characters: every retained excerpt is pure ASCII, so the delivered engine, XeLaTeX with its default Latin Modern text font, reports no missing character.
\begin{align}
\label{algebra-structure}
\iota:H^{*}(BT^{n+1}) \to H^{*}(\mathcal{GQ}_{2n})
\end{align}

\begin{lemma}
\label{surjectivity}
The homomorphism 
$\psi:\mathbb{Z}[\mathcal{GQ}_{2n}]\to H^{*}(\mathcal{GQ}_{2n})$
is surjective.
\end{lemma}

\begin{align}
\label{linear-relation}
f=g_{1}+g_{2}M_{1}+\cdots &+g_{n+1}M_{1}\cdots M_{n} \\
&+g_{n+2}\Delta_{\{n+2,\ldots, 2n+2\}}+g_{n+3}\Delta_{\{n+3,\ldots, 2n+2\}}+\cdots 
+g_{2n+2}\Delta_{\{2n+2\}} \nonumber
\end{align}

\begin{lemma}
\label{injectivity}
The homomorphism 
$\psi:\mathbb{Z}[\mathcal{GQ}_{2n}]\to H^{*}(\mathcal{GQ}_{2n})$
is injective.
\end{lemma}

\begin{proof}[Proof of Lemma~\ref{injectivity}]
It is enough to prove that the following composition map $\varphi$ is injective:
\begin{align*}
\varphi:\mathbb{Z}[\mathcal{GQ}_{2n}]\stackrel{\psi}{\longrightarrow} H^{*}(\mathcal{GQ}_{2n})\hookrightarrow \bigoplus_{v\in V_{2n}}H^{*}(BT^{n+1}) 
%\simeq\bigoplus_{v\in V_{2n}}\mathbb{Z}[\mathcal{GQ}_{2n}]/\langle G_{J}\ |\ v\not\in J \rangle
\simeq \bigoplus_{v\in V_{2n}}\mathbb{Z}[M_{i}\ |\ i\in I_{v}].
\end{align*}
Assume that $\varphi(f)=0$ for an element $f\in \mathbb{Z}[\mathcal{GQ}_{2n}]$. 
We will prove that $f=0$.
In the proof, we use the following restriction map for $w\in V_{2n}$: 
\begin{align*}
\rho_{w}:%\bigoplus_{v\in V_{2n}}\mathbb{Z}[\mathcal{GQ}_{2n}]/\langle G_{J}\ |\ v\not\in J \rangle
\bigoplus_{v\in V_{2n}}\mathbb{Z}[M_{i}\ |\ i\in I_{v}]
\to \mathbb{Z}[M_{i}\ |\ i\in I_{w}]
\end{align*}
and the image of $f\in \mathbb{Z}[\mathcal{GQ}_{2n}]$ by the composition map $\rho_{w}\circ \varphi$ by $f(w)(:=\rho_{w}\circ \varphi(f))$.
The assumption $\varphi(f)=0$ is equivalent to that $\rho_{v}\circ \varphi(f)=f(v)=0\in \mathbb{Z}[M_{i}\ |\ i\in I_{v}]$ for all $v\in V_{2n}$.

In the proof of Lemma~\ref{surjectivity}, especially in Equation~\eqref{linear-relation}, we also show the following fact: for any element $f\in \mathbb{Z}[\mathcal{GQ}_{2n}]$, there exists $g_{i},g_{n}'\in %H^{\deg f-2i}(BT^{n+1})\subset 
\mathbb{Z}[M_{2}-M_{1},\ldots, M_{n+2}-M_{1}] \subset \mathbb{Z}[M_{1},\ldots, M_{n+2}]$ for $i=1,\ldots, 2n$ and $g_{0}\in \mathbb{Z}[M_{2},\ldots, M_{n+2}]$ such that  
\begin{align}
\label{induction_injective}
f&=g_{0}+g_{1}M_{1}+\cdots +g_{n}M_{1}\cdots M_{n}+g_{n}'\Delta_{\{n+2,\ldots, 2n+2\}}+g_{n+1}\Delta_{\{n+3,\ldots, 2n+2\}}+\cdots 
+g_{2n}\Delta_{\{2n+2\}} \\
&=g_{0}+\sum_{i=1}^{n}{g_{i}}M_{1}\cdots M_{i}+X(\Delta), \nonumber
\end{align}
where $X(\Delta)$ is the $\Delta_{K}$ terms.
Note that $\psi(g_{j}), \psi(g_{n}')\in \mathbb{Z}[x_{1},\ldots, x_{n+1}]$ (see \eqref{algebra-structure}) for all $j=0,\ldots, 2n$.
This implies that if there is a vertex $v\in V_{2n}$ such that $g_{j}(v)=0$ (resp.~$g_{n}'(v)=0$), then $g_{j}=0$ (resp.~$g_{n}'=0$).

We first claim that $f$ can be written by $\Delta_{K}$ terms only.
Since $M_{1}(1)=0$ and $X(\Delta)(1)=0$, 
by \eqref{induction_injective}, we have that 
\begin{align*}
g_{0}(1)=f(1)-\left(\sum_{i=1}^{n}{g_{i}}M_{1}\cdots M_{i}+X(\Delta)\right)(1)=f(1)=0.
\end{align*}
%On the other hand, $g_{0}\in \mathbb{Z}[M_{2},\ldots, M_{n+2}]\simeq \mathbb{Z}[M_{2}(1),\ldots, M_{n+2}(1)]\simeq \mathbb{Z}[x_{1},\ldots,x_{n+1}]$. 
%Therefore, we may regard $g_{0}=g_{0}(1)\in \mathbb{Z}[M_{2}(1),\ldots, M_{n+2}(1)]\simeq \mathbb{Z}[x_{1},\ldots,x_{n+1}]$
Therefore, we have $g_{0}=0$.
Similarly, by using $g_{0}=0$ and \eqref{induction_injective}, we have that 
\begin{align*}
g_{1}(2)M_{1}(2)=f(2)-\left(\sum_{i=2}^{n}{g_{i}}M_{1}\cdots M_{i}+X(\Delta)\right)(2)=0.
\end{align*}
Now $g_{1}(2), M_{1}(2)\in \mathbb{Z}[M_{i}\ |\ i\in I_{2}=\{1,3,\ldots,n+2\}]$ and $M_{1}(2)\not=0$.
Since the polynomial ring $\mathbb{Z}[M_{i}\ |\ i\in I_{2}]$ is an integral domain, we see that $g_{1}(2)=0$; therefore, $g_{1}=0$.
%Moreover, because 
%$g_{1}\in \mathbb{Z}[M_{2}-M_{1},\ldots, M_{n+2}-M_{1}]\simeq \mathbb{Z}[x_{1},\ldots, x_{n+1}]$, 
%we have that $g_{1}(2)=g_{1}=0$.
%$g_{1}=h_{1}M_{2}$ for some $h_{1}\in \mathbb{Z}[M_{1},\ldots, M_{n+2}]$.
%Similarly, we have that 
%\begin{align*}
%(h_{1}+g_{2})M_{1}M_{2}(3)=f(3)-\left(\sum_{i=3}^{n}{g_{i}}M_{1}\cdots M_{i}+X(\Delta)\right)(3)=0.
%\end{align*}
%Since $M_{1}M_{2}(3)\not=0$ and $h_{1}+g_{2}\in \mathbb{Z}[M_{1},\ldots, M_{n+2}]$, this implies that $h_{1}M_{1}M_{2}+g_{2}M_{1}M_{2}=h_{2}M_{1}M_{2}M_{3}$ for some $h_{2}\in \mathbb{Z}[M_{1},\ldots, M_{n+2}]$.
Iterating the similar arguments for $i=3,\ldots, n-2$, 
%we can prove that there exists a polynomial $h_{n-1}\in \mathbb{Z}[M_{1},\ldots, M_{n+2}]$ such that 
%\begin{align*}
%h_{n-1}\cdot M_{1}\cdots M_{n}=f-g_{n}M_{1}\cdots M_{n}-X(\Delta).
%\end{align*}
%Put $h=h_{n-1}+g_{n}$.
we also have that $g_{2}=\cdots =g_{n-1}=0$, i.e., 
\begin{align*}
f&=g_{n}M_{1}\cdots M_{n}+X(\Delta) \\
&=g_{n}(\Delta_{\{n+2,\ldots, 2n+2\}}+\Delta_{\{n+1,n+3,\ldots, 2n+2\}})+X(\Delta) \quad ({\rm by\ Relation~3}). % \\
%&=g_{n}(\Delta_{\{n+2,\ldots, 2n+2\}}+\Delta_{\{n+1,n+3,\ldots, 2n+2\}})+g_{n}'\Delta_{\{n+2,\ldots, 2n+2\}}+g_{n+1}\Delta_{\{n+3,\ldots, 2n+2\}}+\cdots +g_{2n}\Delta_{\{2n+2\}}.
\end{align*}
Therefore, if $f(v)=0$ for every $v\in V_{2n}$, then $f$ can be written by the $\Delta_{K}$ terms only; more precisely, 
\begin{align}
\label{middle-eq}
f=g_{n}\Delta_{\{n+1,n+3,\ldots, 2n+2\}}+(g_{n}+g_{n}')\Delta_{\{n+2,\ldots, 2n+2\}}+g_{n+1}\Delta_{\{n+3,\ldots, 2n+2\}}+\cdots 
+g_{2n}\Delta_{\{2n+2\}}.
\end{align}

We next claim that $f=0$ if $f(v)=0$ for every $v\in V_{2n}$.
The equality \eqref{middle-eq} implies that for the vertex $n+1\in V_{2n}$,
\begin{align*}
&g_{n}(n+1)\Delta_{\{n+1,n+3,\ldots, 2n+2\}}(n+1) \\
&=f(n+1)-\left((g_{n}+g_{n}')\Delta_{\{n+2,\ldots, 2n+2\}}+g_{n+1}\Delta_{\{n+3,\ldots, 2n+2\}}+\cdots 
+g_{2n}\Delta_{\{2n+2\}}\right)(n+1)=0.
\end{align*}
Since $\Delta_{\{n+1,n+3,\ldots, 2n+2\}}(n+1)\not=0$, by the similar reason as above, we have $g_{n}=0$.
%$h\in \mathbb{Z}[M_{1},\ldots, M_{n+2}]$, there exists $k_{1}\in \mathbb{Z}[M_{1},\ldots, M_{n+2}]$ such that $h=k_{1}M_{n+1}$.
%So we have 
%\begin{align*}
%f=g_{n}'\Delta_{\{n+2,\ldots, 2n+2\}}+g_{n+1}\Delta_{\{n+3,\ldots, 2n+2\}}+\cdots 
%+g_{2n}\Delta_{\{2n+2\}}.
%\end{align*}
Iterating the similar arguments for $i=n+2,\ldots, 2n+2$, 
we have that $g_{n}'=g_{n+1}=\cdots =g_{2n}=0$.
%By Relation 4, we obtain that $k_{1}M_{n+1}\Delta_{\{n+1,n+3,\ldots, 2n+2\}}=k_{1}\Delta_{\{n+3,\ldots, 2n+2\}}$; therefore, by substituting $h=k_{1}M_{n+1}$ into \eqref{middle-eq}, we have 
%\begin{align*}
%f&=k_{1}\Delta_{\{n+3,\ldots, 2n+2\}}+(k_{1}M_{n+1}+g_{n}')\Delta_{\{n+2,\ldots, 2n+2\}}+g_{n+1}\Delta_{\{n+3,\ldots, 2n+2\}}+\cdots +g_{2n}\Delta_{\{2n+2\}} \\
%&=g_{n}'\Delta_{\{n+2,\ldots, 2n+2\}}+(k_{1}+g_{n+1})\Delta_{\{n+3,\ldots, 2n+2\}}+\cdots +g_{2n}\Delta_{\{2n+2\}}\ ({\rm by\ Relation~1}).
%\end{align*}
%Again we have that $\Delta_{\{n+2,\ldots, 2n+2\}}(n+2)\not=0$ but the other terms are $0$ if they restrict to the vertex $n+2\in V_{2n}$.
%This shows that $g_{n}'=k_{2}M_{n+2}$ for some $k_{2}\in \mathbb{Z}[M_{1},\ldots, M_{n+2}]$. 
%Hence, we have 
%\begin{align*}
%f&=k_{2}M_{n+2}\Delta_{\{n+2,\ldots, 2n+2\}}+(k_{1}+g_{n+1})\Delta_{\{n+3,\ldots, 2n+2\}}+\cdots +g_{2n}\Delta_{\{2n+2\}}\\
%&=(k_{1}+k_{2}+g_{n+1})\Delta_{\{n+3,\ldots, 2n+2\}}+g_{n+2}\Delta_{\{n+4,\ldots, 2n+2\}}+\cdots +g_{2n}\Delta_{\{2n+2\}}\ ({\rm by\ Relation~4}) \\
%&=g'_{n+1}\Delta_{\{n+3,\ldots, 2n+2\}}+g_{n+2}\Delta_{\{n+4,\ldots, 2n+2\}}+\cdots +g_{2n}\Delta_{\{2n+2\}},
%\end{align*}
%where we put $g'_{n+1}=k_{1}+k_{2}+g_{n+1}\in \mathbb{Z}[M_{1},\ldots ,M_{n+2}]$.
%Similarly, we have that 
%\begin{align*}
%& g'_{n+1}\Delta_{\{n+3,\ldots, 2n+2\}}(n+3) \\
%&=f(n+3)-\left(g_{n+2}\Delta_{\{n+4,\ldots, 2n+2\}}+\cdots +g_{2n}\Delta_{\{2n+2\}}\right)(n+3)=0.
%\end{align*}
%Since $\Delta_{\{n+3,\ldots, 2n+2\}}(n+3)\not=0$, we have that 
%\begin{align*}
%g_{n+1}'(n+3)=0\in \mathbb{Z}[M_{i}\ |\ i\in I_{n+3}]=\mathbb{Z}[M_{1},\ldots, M_{n-1},M_{n+1},M_{n+2}].
%\end{align*}
%Recall that by \eqref{algebra-structure} we may regard $g_{n+1}'$ as the element $\iota(g_{n+1}')\in H^{*}(\mathcal{GQ}_{2n})$ such that $\iota(g_{n+1}')(v)=g_{n+1}'$ for every $v\in V_{2n}$.
%Therefore, $g_{n+1}'(n+3)=0$ shows that $g_{n+1}'(v)=0$ for all $v\in V_{2n}$, i.e., $g'_{n+1}=0$.
%%By Relation~2 and $g'_{n+1}\in \mathbb{Z}[M_{1},\ldots ,M_{n+2}]$, this shows that 
%%\begin{align*}
%%g'_{n+1}(=k_{3}M_{n+3})=k_{3}(M_{n+1}+M_{n+2}-M_{n}) 
%%\end{align*}
%%for some 
%%$k_{3}\in \mathbb{Z}[M_{1},\ldots ,M_{n+2}]$.
%Hence, we have 
%\begin{align*}
%f 
%%&=k_{3}(M_{n+1}+M_{n+2}-M_{n})\Delta_{\{n+3,\ldots, 2n+2\}}+g_{n+2}\Delta_{\{n+4,\ldots, 2n+2\}}+\cdots +g_{2n}\Delta_{\{2n+2\}} \\
%&=g_{n+2}\Delta_{\{n+4,\ldots, 2n+2\}}+\cdots +g_{2n}\Delta_{\{2n+2\}}. %({\rm by\ Relation~1}).
%\end{align*}
%By applying the similar arguments for $i=n+4,\ldots, 2n+2$, 
%wesee that 
%$g_{n+2}=\cdots =g_{2n}=0$.
%%$g_{n+2}\Delta_{\{n+4,\ldots, 2n+2\}}=\cdots =g_{2n}\Delta_{\{2n+2\}}=0$.
This establishes that $f=0$.
Consequently, $\varphi$ is injective.
\end{proof}

\end{document}
